%!TEX program = xelatex

\documentclass{article}
% \usepackage[utf8]{inputenc}
\usepackage[UTF8]{ctex}
\usepackage{indentfirst}
\usepackage{lmodern}% http://ctan.org/pkg/lm
\usepackage{setspace}
\usepackage{verbatim}
\usepackage{amsmath}%数学
\usepackage{amsthm}%数学定理、证明
\usepackage{graphicx}%graphicx基于graphics,更方便
% \usepackage{subfig}%图像
\usepackage{subfigure}
\usepackage{booktabs}%表格
\usepackage{tabularx}%表格
\usepackage{multirow}
\usepackage{multicol}
\usepackage{listings}
\usepackage{relsize}
\usepackage[usenames]{color}   
\usepackage{fontspec}

%\usepackage{colortbl}%表格颜色
\usepackage[table]{xcolor}
\usepackage{array}
\usepackage[rightcaption]{sidecap}
\usepackage{hyperref}
\usepackage{float}
\hypersetup{
colorlinks=true,
linkcolor=black
}%解决目录红框问题
\lstset{
         language = vhdl, numbers=left, 
         numberstyle=\tiny,keywordstyle=\color{blue!70},
         commentstyle=\color{red!50!green!50!blue!50},frame=shadowbox,
         rulesepcolor=\color{red!20!green!20!blue!20},basicstyle=\ttfamily
}
\newcommand{\myd}{\;\mathrm{d}}%自定义积分变量命令
% \renewcommand{\listtablename}{表格}%改变listoftable的名字
\renewcommand{\arraystretch}{1.8}%The height of each row, relative to its default height
% \setmonofont{Consolas}
\renewcommand\refname{参考网站}
% \setlength{\tabcolsep}{18pt}%The space between the text and the left/right border of its containing cell
% \setlength{\arrayrulewidth}{0.4mm}%sets the thickness of the borders of the table


\setlength{\parindent}{2em}%控制首行缩进
\addtolength{\parskip}{3pt}%\parskip 控制段落(paragraph)距离
%\singlespacing%单倍行距
%\doublespacing%双倍行距
%\setstretch{1.25}%设置行距为1.25
\onehalfspacing %1.5倍行距
%\begin{doublespacing} 行距环境 \end{doublespacing}
\graphicspath{{./figures/}} % 指定图片所在文件夹
%\newtheorem{definition}[定义]{section}
%定制环境: {环境名}[编号延续]{显示名}[编号层次]

\begin{document}
\begin{titlepage}
        \vspace*{-2.5cm}
	
	\begin{figure}[h]
		\centering
		\includegraphics[width=0.7\linewidth]{zjdx}
	\end{figure}

	\begin{figure}[h]
		\centering
		\includegraphics[width=0.5\linewidth]{QSY}
	\end{figure}
	\vspace{-0.5cm}
	\begin{center}
		\Huge{\textbf{电路与模拟电子技术实验}}\\
		
		\Huge{\textbf{实验报告}}
	\end{center}
	
	
	\vspace*{1.5cm}

    \begin{center}
    \Large
            实验名称\ \ \underline{\makebox[220pt]{直流电压、电流和电阻的测量}} \\ 
            \vspace{0.3cm}
            实验地点\ \ \underline{\makebox[220pt]{东三201-B1实验台}}\\
            \vspace{0.3cm}
            \quad\ 姓\; 名 \ \ \underline{\makebox[220pt]{XXX}}\\
            \vspace{0.3cm}
            \quad\ 学\; 号\ \ \underline{\makebox[220pt]{3100100000}}\\
            \vspace{0.3cm}
            实验日期\ \ \underline{\makebox[220pt]{2020年 9 月 18 日}}\\
            \vspace{0.3cm}
            指导老师\ \ \underline{\makebox[220pt]{XXX}}\\
            \vspace{0.3cm}
            \quad\ 成\; 绩\ \ \underline{\makebox[220pt]{}}\\

              
            

             
    \end{center}
        
    
\end{titlepage}

\newpage
% \setcounter{tocdepth}{5}
\tableofcontents
\thispagestyle{empty}%Removes the page numbering.
% \listoffigures%生成插图
% \listoftables%生成表格
\thispagestyle{empty}%Removes the page numbering.

\newpage
\pagenumbering{arabic}%重新开始标号,阿拉伯数字形式

\section{实验目的和要求}
\section{实验内容和原理}
\section{主要仪器设备}
\section{操作方法和实验步骤}
\section{实验数据记录和处理}
\section{实验结果与分析}
\section{讨论、心得}

\end{document}



%好看的颜色 [HTML] 淡蓝色 CBF5FB
%>{\columncolor[HTML]{AAACED}} 列背景颜色
%\rowcolor, \cellcolor 行背景、单元背景颜色
